%% This file contains a number of tweaks that are typically applied to the main document.
%% They are not enabled by default, but can be enabled by uncommenting the relevant lines.


%%
%% Official support for 2-column teaser figure
%%
\usepackage{cuted} %< for strip environment

%%
%% For auto-referencing of figures (see fig/teaser.tex and \label{} therein)
%%
\usepackage{currfile} %< for \currfilebase macro

%%
%% caption and (global) spacing overrides
%%
\usepackage{caption} %< for "\captionof" commands (teaser)
% \setlength{\abovecaptionskip}{.5em} % space above caption
% \setlength{\belowcaptionskip}{1em} % space below caption

%%
%% Packages used by the paper (graphicx, booktabs, amsmath, xcolor, subcaption are loaded by cvpr.sty)
%%
\usepackage{colortbl}      %< \rowcolor in the tables
\usepackage{algorithm}
\usepackage{algpseudocode}

%%
%% Inline annotations; for predefined colors, refer to "dvipsnames" in the xcolor package:
%% https://tinyurl.com/overleaf-colors
%%
\newcommand{\red}[1]{{\color{red}#1}}
\newcommand{\todo}[1]{\textcolor{red}{TODO: #1}}   %< as in the workshop paper
\newcommand{\TODO}[1]{\textbf{\color{red}[TODO: #1]}}
\newcommand{\mynote}[1]{\textcolor{blue}{Note: #1}}
% [Claude] 2026-11-05: per-author tags and a placeholder tag (user request); dvipsnames colours, distinct from the ones above
\newcommand{\placeholder}[1]{\textcolor{Gray}{[Placeholder: #1]}}
\newcommand{\lulu}[1]{\textcolor{ForestGreen}{[Lulu: #1]}}
\newcommand{\lily}[1]{\textcolor{Magenta}{[Lily: #1]}}
\newcommand{\alec}[1]{\textcolor{Cerulean}{[Alec: #1]}}
\definecolor{pink}{HTML}{D9786A}                    %< "MoMo" in the title

%%
%% Content written by Claude (AI assistant). Everything Claude wrote is set in claudecolor and opens with a
%% "[Claude draft]" tag, so it can be reviewed and rewritten before submission:
%%   \claude{...}                              inline text written by Claude
%%   \begin{claudedraft} ... \end{claudedraft}  a block written by Claude (tagged at its start)
%%   \claudenote{...}                          a note from Claude to the authors (not paper text)
%%
\definecolor{claudecolor}{HTML}{6A3D9A}
\newcommand{\claudetag}{\textsf{\scriptsize[Claude draft]}}
\newcommand{\claude}[1]{{\color{claudecolor}#1}}
\newcommand{\claudenote}[1]{{\color{claudecolor}\textsf{\footnotesize[Claude note: #1]}}}
\newenvironment{claudedraft}{\par\color{claudecolor}\noindent\claudetag\ \ignorespaces}{\par}
%%
%% Disable for camera ready / submission by uncommenting these lines:
%%
% \renewcommand{\TODO}[1]{}
% \renewcommand{\todo}[1]{#1}
% \renewcommand{\lulu}[1]{} \renewcommand{\lily}[1]{} \renewcommand{\alec}[1]{}   % [Claude] 2026-11-05: hide the author tags
% \renewcommand{\claudenote}[1]{}
% \renewcommand{\claudetag}{}
% \definecolor{claudecolor}{HTML}{000000}

%%
%% Work harder in optimizing text layout. Typically shrinks text by 1/6 of page, enable
%% it at the very end of the writing process, when you are just above the page limit.
%%
% \usepackage{microtype}

%%
%% Fine-tune paragraph spacing.
%%
% \renewcommand{\paragraph}[1]{\vspace{.5em}\noindent\textbf{#1.}}

%%
%% Globally adjusts space between figure and caption.
%%
% \setlength{\abovecaptionskip}{.5em}

%%
%% Allows "the use of \paper to refer to the project name"
%% with automatic management of space at the end of the word.
%%
% \usepackage{xspace}
% \newcommand{\paper}{ProjectName\xspace}

%%
%% Commonly used math definitions.
%%
% \DeclareMathOperator*{\argmin}{arg\,min}
% \DeclareMathOperator*{\argmax}{arg\,max}

%%
%% Tighten underline.
%%
% \usepackage{soul}
% \setuldepth{foobar}

%%
%% Reconfigure paragraph (tweaks spacing)
%%
